th are optimistic
reweighting views, but one controls remoteness under a finite negative
measure and the other selects a quality region. Neither problem implies
the other.

\subsection{TOPR Identities and Conditional-versus-Population Constants}
\label{app:topr_relation}

For total surprisal and \(\beta=1\), substituting
\(D_\pi=-\log\pi\) and \(\tau=-\log\mu\) gives
\[
e^{-[D_\pi-\tau]_+}
=e^{-[-\log(\pi/\mu)]_+}
=\min\{\pi/\mu,1\}.
\]
Replacing both log probabilities by their length-normalized values
raises the ratio to the \(1/L\) power, proving
Equations~\eqref{eq:topr_exact_identity}--\eqref{eq:topr_length_normalized}.

For a conditionally reused far-branch sample, the coefficient contains
\(\pi/\mu\), so the vector-field prefactor contains \(\mu^{-1}\).
Solving the asymptotic scalar equation yields
\(\pi(t)\asymp\mu/(Cqt)\) and
\(\pi(t)/\mu\asymp1/(Cqt)\). In the population field, however, the
outer sampling mass gives
\[
\mu\min\{\pi/\mu,1\}=\min\{\pi,\mu\},
\]
and cancels that denominator on the far branch. For Gaussian atoms the
triangle inequality and the one-atom \(G_i\) calculation in
Appendix~\ref{app:gaussian_drpo_boundedness} give weighted constants
first. A \(\mu_{\min}\) expression follows only after imposing finite
support and replacing the weighted terms by their worst case.

\section{Experimental Details}
\label{app:experimental_details}

\subsection{Unified Continuous Environment C-U1}
\label{app:cu1_environment}

C-U1 is the shared controlled continuous environment used by all
continuous mechanism experiments. A context
$s\in\mathbb R^6$
is sampled from
$\mathcal N(0,I_6)$.
For each seed, we independently sample 4096 training contexts and 4096
test contexts from the same distribution. Test contexts never
participate in optimization and are used only to measure
held-out-context generalization.

\paragraph{State-dependent task geometry.}
Each context determines a two-dimensional positive-support center
$a_+(s)$ and a unit task direction $u(s)$. We define
\begin{equation}
\begin{aligned}
    [a_+(s)]_1
    &=
    0.70
    \tanh
    \bigl(
        0.85s_1
        -
        0.30s_2s_3
        +
        0.20\sin(1.6s_4)
    \bigr), \\
    [a_+(s)]_2
    &=
    0.65
    \tanh
    \bigl(
        -0.50s_2
        +
        0.35\cos(1.1s_5)
        +
        0.22s_1s_6
    \bigr).
\end{aligned}
\label{eq:cu1_positive_center}
\end{equation}
The direction angle is
\begin{equation}
\begin{aligned}
    \varphi(s)
    &=
    1.15
    \tanh
    \bigl(
        0.75s_1
        +
        0.50s_3
        -
        0.30s_6
    \bigr) \\
    &\quad+
    0.30\sin(1.35s_2),
\end{aligned}
\label{eq:cu1_direction_angle}
\end{equation}
with
\begin{equation}
    u(s)
    =
    \begin{bmatrix}
        \cos\varphi(s)\\
        \sin\varphi(s)
    \end{bmatrix},
    \qquad
    v(s)
    =
    \begin{bmatrix}
        -u_2(s)\\
        u_1(s)
    \end{bmatrix}.
    \label{eq:cu1_task_basis}
\end{equation}
The unobserved task optimum and the directionally useful negative anchor
are
\begin{equation}
\begin{aligned}
    a_\star(s)
    &=
    a_+(s)+0.70u(s), \\
    a_-(s)
    &=
    a_+(s)-0.50u(s).
\end{aligned}
\label{eq:cu1_task_points}
\end{equation}

\paragraph{Positive and negative actions.}
Each context is associated with four positive actions and eight negative
actions. The positive actions lie on a reward contour of radius
$r_+=0.75$
centered at $a_\star(s)$:
\begin{equation}
\begin{aligned}
    a_{\mathrm{pos},k}(s)
    =
    a_\star(s)
    +
    r_+
    \bigl(
        \cos\phi_k\,u(s)
        +
        \sin\phi_k\,v(s)
    \bigr),
\end{aligned}
\label{eq:cu1_positive_actions}
\end{equation}
where
\begin{equation}
\begin{aligned}
    \{\phi_k\}_{k=1}^{4}
    &=
    \{
        \pi-\theta_1,\,
        \pi+\theta_1,\,
        \pi-\theta_2,\,
        \pi+\theta_2
    \}, \\
    \theta_1
    &=
    0.20,
    \qquad
    \cos\theta_2
    =
    \frac{2(0.70)}{0.75}
    -
    \cos\theta_1.
\end{aligned}
\label{eq:cu1_positive_angles}
\end{equation}
This construction makes the four actions equal-reward while ensuring
that their centroid is exactly $a_+(s)$. It also leaves nonzero
conditional action spread, preventing deterministic maximum-likelihood
variance contraction from being confused with the far-field mechanism.

The eight negative actions lie on a second reward contour of radius
$r_-=1.20$:
\begin{equation}
\begin{aligned}
    a_{\mathrm{neg},j}(s)
    =
    a_\star(s)
    +
    r_-
    \bigl(
        \cos\psi_j\,u(s)
        +
        \sin\psi_j\,v(s)
    \bigr),
\end{aligned}
\label{eq:cu1_negative_actions}
\end{equation}
with
\begin{equation}
\begin{aligned}
    \{\psi_j\}_{j=1}^{8}
    =
    \left\{
        \pi,\,
        \frac{3\pi}{4},\,
        \frac{\pi}{2},\,
        \frac{\pi}{4},\,
        0,\,
        -\frac{\pi}{4},\,
        -\frac{\pi}{2},\,
        -\frac{3\pi}{4}
    \right\}.
\end{aligned}
\label{eq:cu1_negative_angles}
\end{equation}
The first negative action equals $a_-(s)$. All eight negative actions
have the same distance to $a_\star(s)$ and therefore the same reward and
fixed advantage, while having different distances from the current
policy.

\paragraph{Reward and fixed advantages.}
The ground-truth reward is
\begin{equation}
    R(s,a)
    =
    \exp
    \left(
        -
        \frac{
            \lVert a-a_\star(s)\rVert_2^2
        }{
            2(0.75)^2
        }
    \right).
    \label{eq:cu1_reward}
\end{equation}
Advantages are computed once using the fixed baseline $0.40$:
\begin{equation}
    \widehat A(s,a)
    =
    R(s,a)-0.40,
    \label{eq:cu1_advantage}
\end{equation}
and remain frozen throughout optimization. Consequently, all four
positive actions have positive advantages, all eight negative actions
have negative advantages, and negative-action quality is exactly matched
within each context.

\paragraph{Policy parameterization.}
The policy is an isotropic state-conditioned Gaussian,
\begin{equation}
    \pi_\theta(a\mid s)
    =
    \mathcal N
    \left(
        \mu_\theta(s),
        \sigma_\theta^2(s)I_2
    \right).
    \label{eq:cu1_policy}
\end{equation}
Both $\mu_\theta(s)$ and the scalar
$\log\sigma_\theta(s)$
are produced by a shared two-layer ReLU MLP with 64 hidden units per
layer and separate output heads. The initial standard deviation is
$0.60$, and the learnable-variance experiments use no artificial
variance clamp.  The continuous mechanism results reported in the main
text use the fixed-variance branch unless explicitly labeled otherwise. A finite crossing of
$\log\sigma_\theta(s)<-12$
is recorded as a support or variance-boundary event, whereas nonfinite
parameters or outputs are reported separately as NaN/Inf numerical
failure.

Training minibatches are sampled by context, and all actions associated
with a selected context are retrieved together. Thus, the replicated
actions within one context are not treated as independent contexts.
Positive and negative objectives are averaged separately before their
relative coefficient is applied, preventing the four-to-eight action
count ratio from mechanically changing the total negative-update mass.

\paragraph{Shared use across continuous protocols.}
All C-U1 experiments reuse the same contexts, reward geometry, fixed
advantages, and policy architecture. E1 compares equal-advantage
negative actions at different learner-relative distances. E2 applies
only positive updates and measures the same negative actions as phantom
probes. E3 dynamically partitions negative actions using the
standardized distance
\begin{equation}
    d_\theta(s,a)
    =
    \frac{
        \lVert a-\mu_\theta(s)\rVert_2
    }{
        \sigma_\theta(s)
    },
    \label{eq:cu1_standardized_distance}
\end{equation}
with the prespecified near--far threshold $d_\theta=5.0$, and applies
targeted near- and far-field interventions. E4 uses the aligned action
$a_-(s)$ as directionally useful local negative feedback and the
remaining contour actions as additional far-field pressure. The
experiment-specific coefficients, training horizons, seeds, taper
calibrations, and terminal criteria are reported separately with their
corresponding protocols.

\subsection{Controlled Categorical Environment D-U1}
\label{app:du1_environment}

D-U1 denotes the controlled categorical environment family used to
study persistent negative suppression, probability-support boundaries,
and shared-semantic generalization. Its E6 instantiation uses a fixed
catalogue of 64 unordered actions, each associated with a
four-dimensional unit semantic vector. Action identifiers are randomly
permuted for every seed, so numerical action indices carry no geometric
ordering.

For each seed, we independently sample 2048 training contexts and 2048
test contexts from
$s\sim\mathcal N(0,I_6)$.
The two splits follow the same context distribution, and test contexts
never participate in optimization. Results on this split are therefore
reported as held-out-context or unseen-state generalization rather than
OOD generalization.

\paragraph{Semantic action catalogue.}
Let
$\{e_i\}_{i=1}^{K}\subset\mathbb R^4$,
with $K=64$, denote the unit reward embeddings of the categorical
actions. They are sampled once per seed and randomly assigned to action
identifiers. For each context, two fixed linear maps generate a
demonstrated semantic target and a candidate improvement direction:
\begin{equation}
\begin{aligned}
    t_+(s)
    &=
    \operatorname{unit}
    \!\left(
        W_+^\top s
    \right), \\
    \widetilde u(s)
    &=
    W_u^\top s
    -
    \left\langle
        W_u^\top s,
        t_+(s)
    \right\rangle
    t_+(s), \\
    u(s)
    &=
    \operatorname{unit}
    \!\left(
        \widetilde u(s)
    \right).
\end{aligned}
\label{eq:du1_semantic_basis}
\end{equation}
The hidden optimal target and the local negative target are
\begin{equation}
\begin{aligned}
    t_\star(s)
    &=
    \operatorname{unit}
    \!\left(
        t_+(s)+\delta u(s)
    \right), \\
    t_-(s)
    &=
    \operatorname{unit}
    \!\left(
        t_+(s)-\delta u(s)
    \right),
\end{aligned}
\qquad
\delta=0.45.
\label{eq:du1_semantic_targets}
\end{equation}

\paragraph{Reward and action roles.}
The ground-truth reward of action $i$ is determined by its semantic
alignment with the hidden target:
\begin{equation}
    R(s,i)
    =
    \frac{1}{2}
    \left(
        1+
        t_\star(s)^\top e_i
    \right).
    \label{eq:du1_reward}
\end{equation}
The hidden optimal action is
\begin{equation}
    i_\star(s)
    =
    \arg\max_i
    t_\star(s)^\top e_i.
    \label{eq:du1_hidden_optimum}
\end{equation}
It is excluded from all positive demonstrations.

For each context, the four positive actions are the available actions
with the largest similarity to $t_+(s)$. After excluding the hidden
optimum and positive actions, the local negative action is the action
most aligned with $t_-(s)$. The four far-negative actions are selected
by their alignment with $-t_+(s)$. Thus,
\begin{equation}
\begin{aligned}
    \mathcal P(s)
    &=
    \operatorname{Top4}_{i\neq i_\star(s)}
    \,
    t_+(s)^\top e_i, \\
    i_{\mathrm{local}}(s)
    &=
    \arg\max_{i\notin
        \{i_\star(s)\}\cup\mathcal P(s)}
    t_-(s)^\top e_i, \\
    \mathcal F(s)
    &=
    \operatorname{Top4}_{i\notin
        \{i_\star(s),i_{\mathrm{local}}(s)\}
        \cup\mathcal P(s)}
    \,
    \bigl[-t_+(s)^\top e_i\bigr].
\end{aligned}
\label{eq:du1_action_roles}
\end{equation}
The positive and negative advantages are fixed throughout training:
\begin{equation}
    \widehat A(s,i)
    =
    \begin{cases}
        +1,
        & i\in\mathcal P(s),\\
        -1,
        & i=i_{\mathrm{local}}(s)
          \text{ or }i\in\mathcal F(s).
    \end{cases}
    \label{eq:du1_fixed_advantages}
\end{equation}
Consequently, differences between local and far-negative updates cannot
be attributed to different advantage magnitudes.

\paragraph{Shared-semantic categorical policy.}
The policy uses a two-layer tanh MLP with 64 hidden units per layer to
map the context to a unit semantic direction
$q_\theta(s)\in\mathbb R^4$.
Let $\widetilde e_i$ denote the action embedding used by the policy. The
categorical logits and probabilities are
\begin{equation}
\begin{aligned}
    \ell_{\theta,i}(s)
    &=
    \kappa_\theta(s)
    q_\theta(s)^\top
    \widetilde e_i, \\
    \pi_\theta(i\mid s)
    &=
    \frac{
        \exp\!\left(\ell_{\theta,i}(s)\right)
    }{
        \sum_{j=1}^{K}
        \exp\!\left(\ell_{\theta,j}(s)\right)
    }.
\end{aligned}
\label{eq:du1_policy}
\end{equation}
The concentration is either fixed at
$\kappa_\theta(s)=8$
or learned as
\begin{equation}
    \kappa_\theta(s)
    =
    \operatorname{softplus}
    \!\left(
        c_\theta(s)
    \right)
    +
    0.05,
    \label{eq:du1_concentration}
\end{equation}
with initial value 8 and no upper clamp.

In the semantically aligned condition,
$\widetilde e_i=e_i$.
In the shuffled control, the same policy embeddings are randomly
permuted across action identifiers while the reward embeddings and
training samples remain unchanged. This intervention preserves the
catalogue size and categorical parameterization but breaks the
correspondence between policy-side similarity and ground-truth semantic
utility.

\paragraph{Evaluation and failure events.}
We evaluate expected semantic reward, probability assigned to the hidden
optimal action, entropy, effective support, and normalized semantic
extrapolation on both training and held-out contexts. Task-performance
collapse, effective-support or concentration-boundary events, and
NaN/Inf numerical failure are reported separately. A policy can
therefore retain high reward while still approaching a support boundary;
the two outcomes are not treated as equivalent.

\paragraph{Shared use across categorical protocols.}
E6-A fixes the concentration and scans the strength of the local
negative update to identify the transition from the positive-only
ceiling to useful semantic extrapolation and then to performance
reversal. E6-B learns the concentration and applies near/far causal
interventions and matched controls. E6-C compares aligned and shuffled
policy embeddings to determine whether improvements on the hidden action
require the specified semantic correspondence.

E5 serves a distinct but complementary role. Its direct-softmax
diagnostic and categorical causal reconstruction isolate persistent
surprisal growth, logit-gap growth, and simplex-boundary suppression
under repeated negative updates. It shares the fixed-advantage and
near/far intervention logic of D-U1, but it does not use the 64-action
shared-semantic catalogue and does not support the E6 semantic
generalization claim.

\subsection{D4RL Hopper Datasets}
\label{app:hopper_datasets}

Hopper is a continuous-control task with an
11-dimensional observation space and a three-dimensional bounded action
space. We use the D4RL
\texttt{hopper-medium},
\texttt{hopper-medium-replay}, and
\texttt{hopper-medium-expert}
datasets. The medium dat